%% Covering Factor Definitions for the LaRT Clumpy Medium
%% Author: Kwang-Il Seon
\documentclass[english,a4paper]{my_memo}
\synctex=1
\usepackage{amsmath}
\usepackage{amssymb}
\usepackage{microtype}
\usepackage{babel}

\begin{document}

\title{Covering Factor Definitions in the LaRT Clumpy Medium}
\author{Kwang-Il Seon}
\date{\docmoddate}
\maketitle

\tableofcontents
\bigskip

%=============================================================================
\section{Setup and Notation}
%=============================================================================

We consider a spherical domain of radius $R$ filled with a population of
spherical clumps.  In the present LaRT implementation each clump has the
same radius $r_c$ and the same internal opacity, but the formulas below are
written so that a generalization to position-dependent
$r_c(\mathbf{r})$ and clump number density $n_c(\mathbf{r})$ is immediate.

The fundamental local quantity is the differential covering rate per unit
path length,
\begin{equation}
   \kappa_{\rm cov}(\mathbf{r}) \;\equiv\; n_c(\mathbf{r})\,\pi\,r_c(\mathbf{r})^2 ,
   \label{eq:kappa_cov}
\end{equation}
which has units of inverse length and represents the expected number of
clumps a sightline encounters per unit length when passing through the
neighborhood of $\mathbf{r}$.

For uniform profiles
$n_c(\mathbf{r})=n_c^{(0)}\equiv N/[(4/3)\pi R^{3}]$ and
$r_c(\mathbf{r})=r_c$, so
$\kappa_{\rm cov} = N\,r_c^{2}/((4/3)R^{3})$.

%=============================================================================
\section{Current LaRT Definition (Uniform Clumps)}
%=============================================================================

LaRT currently uses
\begin{equation}
   \boxed{\;
   f_{\rm cov}^{\rm LaRT} \;=\; \tfrac{3}{4}\,N\,\Bigl(\frac{r_c}{R}\Bigr)^{\!2}\; }
   \label{eq:fcov_LaRT}
\end{equation}
as the covering factor associated with $N$ clumps of radius $r_c$ packed
into a sphere of radius $R$ (see \texttt{clump\_mod.f90}, lines~93 and~136).
The companion volume filling factor is
\begin{equation}
   f_{\rm vol}^{\rm LaRT} \;=\; N\,\Bigl(\frac{r_c}{R}\Bigr)^{\!3}.
   \label{eq:fvol_LaRT}
\end{equation}

Equation~\eqref{eq:fcov_LaRT} is exactly the radial line-of-sight integral
of $\kappa_{\rm cov}$ from the box center to the surface,
\begin{equation}
   f_{\rm cov}^{\rm LaRT}
   \;=\; \int_{0}^{R}\! \kappa_{\rm cov}\,dr
   \;=\; n_c^{(0)}\,\pi\,r_c^{2}\,R
   \;=\; \tfrac{3}{4}\,N\,(r_c/R)^{2},
\end{equation}
i.e.~the average number of clumps a photon emitted at the center crosses on
its way out.  This is the natural Monte-Carlo interpretation: it equals the
expected count of clump intersections along a single radial trajectory
emerging from the source.

%=============================================================================
\section{Generalized Definitions for Non-Uniform Profiles}
%=============================================================================

When $r_c$ and $n_c$ depend on position, the single number $f_{\rm cov}$ must
be defined by a particular averaging prescription.  Three natural choices
are described below.  All three reduce to Eq.~\eqref{eq:fcov_LaRT} (up to a
constant prefactor) in the uniform limit.

\subsection{Radial Line-of-Sight Integral (current LaRT convention)}
\label{sec:los}

Integrate $\kappa_{\rm cov}$ along a single radial sightline from the box
center to the surface,
\begin{equation}
   f_{\rm cov}^{\rm LOS}
   \;=\; \int_{0}^{R}\! n_c(r)\,\pi\,r_c(r)^{2}\,dr.
   \label{eq:fcov_los}
\end{equation}
For a spherically-symmetric profile the result is direction-independent.

\paragraph{Uniform limit.} Substituting $n_c=n_c^{(0)}$ and constant $r_c$
recovers exactly Eq.~\eqref{eq:fcov_LaRT}.

\paragraph{Physical meaning.} The expected number of clumps a photon
launched at the center traverses on its way to the surface.  This is the
quantity that controls the radiative-transfer outcome for a centrally
located source most directly.

\subsection{Impact-Parameter Average (Random External Sightline)}
\label{sec:impact}

A distant external observer sees the box covered by chords of varying
impact parameter $b\in[0,R]$.  Averaging the covering along such chords
with the natural area weight yields
\begin{equation}
   \langle f_{\rm cov}\rangle_{b}
   \;=\; \frac{1}{\pi R^{2}}
         \int_{0}^{R}\! 2\pi b\,f_{\rm cov}(b)\,db,
   \label{eq:fcov_impact}
\end{equation}
where $f_{\rm cov}(b)$ is the line integral of $\kappa_{\rm cov}$ along the
chord at impact parameter $b$.

\paragraph{Uniform limit.} The chord length is $2\sqrt{R^{2}-b^{2}}$, so
\begin{align}
   \langle f_{\rm cov}\rangle_{b}^{\rm uniform}
   &= n_c^{(0)}\,\pi\,r_c^{2}\,
      \frac{4}{R^{2}}\!\int_{0}^{R}\!b\sqrt{R^{2}-b^{2}}\,db \nonumber \\
   &= \tfrac{4}{3}\,n_c^{(0)}\,\pi\,r_c^{2}\,R
   \;=\; N\,\Bigl(\frac{r_c}{R}\Bigr)^{\!2},
   \label{eq:fcov_impact_uniform}
\end{align}
which exceeds Eq.~\eqref{eq:fcov_LaRT} by a factor $4/3$.

\paragraph{Physical meaning.} The covering factor as inferred from
absorption observations of a background source for which all sightlines
through the cloud are equally likely (uniform illumination by a
distant continuum source).

\subsection{Volume-Weighted Average}
\label{sec:volume}

Volume-average the local covering rate $\kappa_{\rm cov}(\mathbf{r})$ and
multiply by the radius to obtain a length-scaled covering,
\begin{equation}
   f_{\rm cov}^{\rm vol}
   \;\equiv\; R\,\bigl\langle\kappa_{\rm cov}\bigr\rangle_{V}
   \;=\; \frac{R}{V_{\rm box}}\!\int n_c(r)\,\pi\,r_c(r)^{2}\,dV
   \;=\; \frac{3\pi}{R^{2}}\!\int_{0}^{R}\!n_c(r)\,r_c(r)^{2}\,r^{2}\,dr.
   \label{eq:fcov_volume}
\end{equation}

\paragraph{Uniform limit.}
$f_{\rm cov}^{\rm vol,\,uniform} = R\,n_c^{(0)}\,\pi\,r_c^{2} =
\tfrac{3}{4}\,N\,(r_c/R)^{2}$, identical to Eq.~\eqref{eq:fcov_LaRT}.

\paragraph{Physical meaning.} A bulk diagnostic of how much projected
clump area exists in the box, normalized by the domain size.  The radial
weight $r^{2}$ heavily favors outer clumps over inner clumps and hence
differs from the LOS integral whenever the profile is non-trivial.

%=============================================================================
\section{Equivalence in the Uniform Limit}
%=============================================================================

In the uniform-clump limit the three definitions are related by simple
constants, summarized below:

\begin{center}
\renewcommand{\arraystretch}{1.4}
\begin{tabular}{@{}lll@{}}
\hline
Definition & Symbol & Uniform-limit value \\
\hline
Radial LOS (current LaRT)        & $f_{\rm cov}^{\rm LOS}$            & $\tfrac{3}{4}N(r_c/R)^{2}$ \\
Impact-parameter average         & $\langle f_{\rm cov}\rangle_{b}$   & $N(r_c/R)^{2}$ \\
Volume-weighted                  & $f_{\rm cov}^{\rm vol}$            & $\tfrac{3}{4}N(r_c/R)^{2}$ \\
\hline
\end{tabular}
\end{center}

The radial-LOS and volume-weighted definitions coincide in the uniform
case but differ in general (the volume weight $4\pi r^{2}$ is absent from
the LOS form).  The impact-parameter average exceeds the LOS form by a
constant factor $4/3$ in the uniform limit.

%=============================================================================
\section{Behavior for Non-Uniform Profiles}
%=============================================================================

The three averages diverge once $r_c$ or $n_c$ depend on position.  Two
illustrative cases:

\paragraph{(a) Thin shell of clumps at $r=r_{0}$.}
Let $n_c(r) = (N/4\pi r_{0}^{2})\,\delta(r-r_{0})$ with constant $r_c$.

\begin{itemize}
\item LOS: $f_{\rm cov}^{\rm LOS} = (N/4\pi r_{0}^{2})\,\pi\,r_c^{2} = N r_c^{2}/(4 r_{0}^{2})$.\\
  Independent of $R$; only the shell radius matters.
\item Volume-weighted: $f_{\rm cov}^{\rm vol} = (3\pi/R^{2})\,(N/4\pi r_{0}^{2})\,r_c^{2}\,r_{0}^{2}\,(\text{shell width factor})$. \\
  Suppressed by $(r_{0}/R)^{0}$ but multiplied by $r^{2}$ inside the integral, so a shell at large $r_0$ contributes more than a shell at small $r_0$.
\item Impact-parameter average: depends on $r_{0}/R$ through the chord-length geometry.
\end{itemize}

\paragraph{(b) Power-law clump density $n_c(r)\propto r^{-\alpha}$.}
With constant $r_c$ and $\alpha<3$ to keep the total finite,
\begin{align}
   f_{\rm cov}^{\rm LOS}    &\propto \int_{0}^{R}\! r^{-\alpha}\,dr   \;\propto\; R^{1-\alpha}, \\
   f_{\rm cov}^{\rm vol}    &\propto \int_{0}^{R}\! r^{-\alpha}r^{2}\,dr \;\propto\; R^{3-\alpha}.
\end{align}
The LOS integral is dominated by inner clumps (for $\alpha>1$), while the
volume-weighted form is dominated by outer clumps for any $\alpha<3$.
This contrast is the key reason the two definitions cannot be used
interchangeably outside the uniform limit.

%=============================================================================
\section{Recommendation for the Generalized Implementation}
%=============================================================================

The radial line-of-sight integral, Eq.~\eqref{eq:fcov_los}, is the
recommended default because:

\begin{enumerate}
\item It coincides with the current LaRT formula
      Eq.~\eqref{eq:fcov_LaRT} in the uniform limit, so existing input
      files retain their meaning.
\item It is the quantity most directly probed by a centrally located
      photon source; for radiative-transfer diagnostics this is a more
      faithful indicator than a volume-weighted bulk number.
\item It is a single one-dimensional integral and is straightforward
      to evaluate numerically in the helper routine
      \texttt{derive\_clump\_\-normalization}.
\end{enumerate}

The impact-parameter average, Eq.~\eqref{eq:fcov_impact}, can optionally
be reported alongside as a diagnostic for comparison with absorption-line
observations of background sources.  The volume-weighted definition,
Eq.~\eqref{eq:fcov_volume}, is included here for completeness but is not
recommended as the primary input quantity because of its strong outer-shell
weighting.

\end{document}
